% ======================================================================
% Chapter 20: Gauge Theories with Spontaneous Symmetry Breaking
% (book pages 689-728; PDF pages 712-753)
% ======================================================================
\chapter{Gauge Theories with Spontaneous Symmetry Breaking}

In the course of this book, we have discussed three distinct fashions in which
symmetries can be realized in a quantum field theory.  The simplest case is a
global symmetry that is manifest, leading to particle multiplets with
restricted interactions.  A second possibility is a global symmetry that is
spontaneously broken.  Then, as discussed in Chapter~11,%
\footnote{Section 11.1 is necessary background for the present chapter, but
the rest of Chapter 11 is not.} the symmetry currents are still conserved and
interactions are similarly restricted, but the vacuum state does not respect
the symmetry and the particles do not form obvious symmetry multiplets.
Instead, such a theory contains massless particles, Goldstone bosons, one for
each generator of the spontaneously broken symmetry.  The third case is that
of a local, or gauge, symmetry.  As we saw in Chapter~15, such a symmetry
requires the existence of a massless vector field for each symmetry generator,
and the interactions among these fields are highly restricted.

It is now only natural to consider a fourth possibility: What happens if we
include both local gauge invariance and spontaneous symmetry breaking in the
same theory?  In this chapter and the next, we will find that this combination
of ingredients leads to new possibilities for the construction of quantum field
theory models.  We will see that spontaneous symmetry breaking requires gauge
vector bosons to acquire mass.  However, the interactions of these massive
bosons are still constrained by the underlying gauge symmetry, and these
constraints can have observable consequences.

In elementary particle physics, the principal application of spontaneously
broken local symmetry is in the currently accepted model of weak
interactions.  This model, due to Glashow, Weinberg, and Salam, is introduced
in Section~20.2.  There we will see that it makes a number of precise and
successful predictions for weak interaction phenomena.  Remarkably, this
model also unifies the weak interactions with electromagnetism in a single
larger gauge theory.

\section{The Higgs Mechanism}

In this section we analyze some simple examples of gauge theories with
spontaneous symmetry breaking.  We begin with an Abelian gauge theory, and
then study several examples of non-Abelian models.

\subsection{An Abelian Example}

As our first example, consider a complex scalar field coupled both to itself
and to an electromagnetic field:
\begin{equation}
  \mathcal{L} = -\frac{1}{4}(F_{\mu\nu})^2 + |D_\mu\phi|^2 - V(\phi),
  \label{eq:20.1}
\end{equation}
with $D_\mu = \partial_\mu + ieA_\mu$.  This Lagrangian is invariant under the
local $U(1)$ transformation
\begin{equation}
  \phi(x) \to e^{i\alpha(x)}\phi(x),\qquad
  A_\mu(x) \to A_\mu(x) - \frac{1}{e}\partial_\mu\alpha(x) .
  \label{eq:20.2}
\end{equation}
If we choose the potential in $\mathcal{L}$ to be of the form
\begin{equation}
  V(\phi) = -\mu^2\phi^{*}\phi + \frac{\lambda}{4}(\phi^{*}\phi)^2 ,
  \label{eq:20.3}
\end{equation}
with $\mu^2 > 0$, the field $\phi$ will acquire a vacuum expectation value and
the $U(1)$ global symmetry will be spontaneously broken.  The minimum of this
potential occurs at
\begin{equation}
  \langle\phi\rangle = \phi_0 = \Bigl( \frac{\mu^2}{\lambda} \Bigr)^{1/2} ,
  \label{eq:20.4}
\end{equation}
or at any other value related by the $U(1)$ symmetry \eqref{eq:20.2}.

Let us expand the Lagrangian \eqref{eq:20.1} about the vacuum state
\eqref{eq:20.4}.  Decompose the complex field $\phi(x)$ as
\begin{equation}
  \phi(x) = \phi_0 + \frac{1}{\sqrt{2}}\bigl( \phi_1(x) + i\phi_2(x) \bigr) .
  \label{eq:20.5}
\end{equation}
The potential \eqref{eq:20.3} is rewritten
\begin{equation}
  V(\phi) = \frac{1}{2}(2\mu^2)\phi_1^2 + \lambda\phi_0\,\phi_1^3
    + \mathcal{O}(\phi^4) ,
  \label{eq:20.6}
\end{equation}
so that the field $\phi_1$ acquires the mass $m = \sqrt{2}\mu$ and $\phi_2$ is
the massless Goldstone boson.  So far, this whole analysis follows that in
Section~11.1.  But now consider how the kinetic energy term of $\phi$ is
transformed.  Inserting the expansion \eqref{eq:20.5}, we rewrite
\begin{equation}
  |D_\mu\phi|^2 = \frac{1}{2}(\partial_\mu\phi_1)^2 + \frac{1}{2}
    (\partial_\mu\phi_2)^2 + \sqrt{2}e\phi_0\cdot A^\mu(\partial_\mu\phi_2)
      + \cdots ,
  \label{eq:20.7}
\end{equation}
where we have omitted terms cubic and quartic in the fields $A^\mu$,
$\phi_1$, and $\phi_2$.  The last term written explicitly in \eqref{eq:20.7}
is a photon mass term
\begin{equation}
  \Delta\mathcal{L} = \frac{1}{2}m_A^2\,A_\mu A^\mu ,
  \label{eq:20.8}
\end{equation}
where the mass
\begin{equation}
  m_A = \sqrt{2}e\phi_0
  \label{eq:20.9}
\end{equation}
arises from the nonvanishing vacuum expectation value of $\phi$.  Notice that
the sign of this mass term is correct; the physical spacelike components of
$A^\mu$ appear in \eqref{eq:20.8} as
\begin{equation}
  \Delta\mathcal{L} = -\frac{1}{2}m_A^2\,\mathbf{A}^2 ,
\end{equation}
with the correct sign for a potential energy term.

In Chapter~7, and again in Chapter~16 for the non-Abelian case, we argued
that a gauge boson cannot obtain a mass, unless this mass term is associated
with a pole in the vacuum polarization amplitude.  There is a counterexample
to this result in two-dimensional spacetime; there, as we saw in Section~19.1,
a pole of the required form can arise from the infrared singularity generated
by a massless fermion pair.  However, in four dimensions, a pole in the vacuum
polarization amplitude can be created only by a massless scalar particle.
Typically, in situations with unbroken symmetry, no such particle is
available.

However, a model with a spontaneously broken continuous symmetry must have
massless Goldstone bosons.  These scalar particles have the quantum numbers
of the symmetry currents, and therefore have just the right quantum numbers
to appear as intermediate states in the vacuum polarization.  In the model we
are now discussing, we can see this pole arise explicitly in the following way:
The third term in Eq.~\eqref{eq:20.7} couples the gauge boson directly to the
Goldstone boson $\phi_2$; this gives a vertex of the form
\begin{equation}
  \text{($\phi_2$-gauge vertex)} = i\sqrt{2}e\phi_0(-ik^\mu)
  = im_A\,k^\mu .
  \label{eq:20.10}
\end{equation}
If we also treat the mass term \eqref{eq:20.8} as a vertex in perturbation
theory, then the leading-order contributions to the vacuum polarization
amplitude give the expression
\begin{equation}
  \Pi^{\mu\nu}(k) = im_A^2\,g^{\mu\nu} + \bigl( m_A k^\mu \bigr)\,
    \frac{i}{k^2}\,\bigl( -m_A k^\nu \bigr)
  = im_A^2\Bigl( g^{\mu\nu} - \frac{k^\mu k^\nu}{k^2} \Bigr) .
  \label{eq:20.11}
\end{equation}
The Goldstone boson supplies exactly the right pole to make the vacuum
polarization amplitude properly transverse.

Although the Goldstone boson plays an important formal role in this theory,
it does not appear as an independent physical particle.  The easiest way to
see this is to make a particular choice of gauge, called the \emph{unitarity
gauge}.  Using the local $U(1)$ gauge symmetry \eqref{eq:20.2}, we can choose
$\alpha(x)$ in such a way that $\phi(x)$ becomes real-valued at every point
$x$.  With this choice, the field $\phi_2$ is removed from the theory.  The
Lagrangian \eqref{eq:20.1} becomes
\begin{equation}
  \mathcal{L} = -\frac{1}{4}(F_{\mu\nu})^2 + (\partial_\mu\phi_1)^2
    + e^2\phi_1^2A_\mu A^\mu - V(\phi) .
  \label{eq:20.12}
\end{equation}
If the potential $V(\phi)$ favors a nonzero vacuum expectation value of
$\phi$, the gauge field acquires a mass; it also retains a coupling to the
remaining, physical field $\phi_1$.

This mechanism, by which spontaneous symmetry breaking generates a mass for
a gauge boson, was explored and generalized to the non-Abelian case by Higgs,
Kibble, Guralnik, Hagen, Brout, and Englert, and is now known as the
\emph{Higgs mechanism}.  However, this mechanism had an earlier application
to the theory of superconductivity.  In Chapter~8, we constructed the Landau
description of a second-order phase transition.  To describe a superconductor,
Landau and Ginzburg coupled this theory to an external electromagnetic field;
they obtained precisely the Lagrangian \eqref{eq:20.1}.  Since the gauge field
acquires a nonzero mass, external electromagnetic fields penetrate a
superconductor only to the depth $m_A^{-1}$.  This explains the Meissner
effect, the observed exclusion of macroscopic magnetic fields from a
superconductor.

The role of the Goldstone boson in the Higgs mechanism is intricate, and
seems mysterious at this level of the discussion.  We first saw that the
involvement of the Goldstone boson is necessary, as a matter of principle, in
order for the gauge boson to acquire a mass.  We then saw that the Goldstone
boson can be formally eliminated from the theory.  However, we might argue
that the Goldstone boson has not completely disappeared.  A massless vector
boson has only two physical polarization states; we saw in Chapter~16 that the
longitudinal polarization state cannot be produced, and appears in the
formalism only to cancel other unphysical contributions.  However, a massive
vector boson must have three physical polarization states: In its rest frame,
it is a spin-1 object, which can make no distinction between transverse and
longitudinal polarizations.  It is tempting to say that the gauge boson
acquired its extra degree of freedom by \emph{eating} the Goldstone boson.
In Sections~21.1 and~21.2 we will clarify this picture, by studying the
quantization and gauge-fixing of spontaneously broken gauge theories.

\subsection{Systematics of the Higgs Mechanism}

The Higgs mechanism extends straightforwardly to systems with non-Abelian
gauge symmetry.  It is not difficult to derive the general relation by which a
set of scalar field vacuum expectation values leads to the appearance of gauge
boson masses.  Let us work out this relation and then apply it in a number of
examples.

Consider a system of scalar fields $\phi_i$ that appear in a Lagrangian
invariant under a symmetry group $G$, represented by the transformation
\begin{equation}
  \phi_i \to \bigl( 1 + i\alpha^a t^a \bigr)_{ij}\phi_j .
  \label{eq:20.13}
\end{equation}
It is convenient to write the $\phi_i$ as real-valued fields, for example,
writing $n$ complex fields as $2n$ real fields.  Then the group
representation matrices $t^a$ must be pure imaginary and, since they are
Hermitian, antisymmetric.  Let us write
\begin{equation}
  t^a_{ij} = iT^a_{ij} ,
  \label{eq:20.14}
\end{equation}
so that the $T^a$ are real and antisymmetric.

If we promote the symmetry group $G$ to a local gauge symmetry, the
covariant derivative on the $\phi_i$ is
\begin{equation}
  D_\mu\phi = \bigl( \partial_\mu - igA^a_\mu t^a \bigr)\phi
  = \bigl( \partial_\mu + gA^a_\mu T^a \bigr)\phi .
  \label{eq:20.15}
\end{equation}
Then the kinetic energy term for the $\phi_i$ is
\begin{equation}
  |D_\mu\phi|^2 = \frac{1}{2}(\partial_\mu\phi_i)^2
    + gA^{a\mu}\,\partial_\mu\phi_i\,(T^a\phi)_i
    + \frac{1}{2}g^2A^a_\mu A^{b\mu}\,(T^a\phi)_i(T^b\phi)_i .
\end{equation}
Now let the $\phi_i$ acquire vacuum expectation values
\begin{equation}
  \langle\phi\rangle = \phi_0 ,
  \label{eq:20.16}
\end{equation}
and expand the $\phi_i$ about these values.  The last term in Eq.~\eqref{eq:20.15}
contains a term with the structure of a gauge boson mass,
\begin{equation}
  \Delta\mathcal{L} = \frac{1}{2}m^2_{ab}A^a_\mu A^{b\mu} ,
  \label{eq:20.17}
\end{equation}
with the mass matrix
\begin{equation}
  m^2_{ab} = g^2\,(T^a\phi_0)_i\,(T^b\phi_0)_i .
  \label{eq:20.18}
\end{equation}
This matrix is positive semidefinite, since any diagonal element, in any
basis, has the form
\begin{equation}
  m^2_{aa} = g^2|T^a\phi_0|^2 \ge 0 \qquad \text{(no sum)} .
\end{equation}
Thus, generically, all of the gauge bosons will receive positive masses.
However, it may be that some particular generator $T^a$ of $G$ leaves the
vacuum invariant:
\begin{equation}
  T^a\phi_0 = 0 .
\end{equation}
In that case, the generator $T^a$ will give no contribution to
\eqref{eq:20.18}, and the corresponding gauge boson will remain massless.

As in the Abelian case, the gauge boson propagator receives a contribution
from the Goldstone bosons, which is necessary to make the vacuum polarization
amplitude transverse.  To compute this contribution, we need the vertex that
mixes gauge bosons and Goldstone bosons.  This comes from the second term of
the Lagrangian \eqref{eq:20.15}.  When we insert the vacuum expectation value
of the scalar field \eqref{eq:20.16}, this term becomes
\begin{equation}
  \Delta\mathcal{L} = gA^{a\mu}\,\partial_\mu\phi_i\,(T^a\phi_0)_i .
  \label{eq:20.19}
\end{equation}
This interaction term does not involve all of the components of $\phi$---only
those that are Goldstone bosons.

As an example of the non-Abelian Higgs mechanism, consider an $SU(2)$ gauge
theory with a scalar field in the fundamental (doublet) representation.  The
vacuum expectation value breaks $SU(2)$ completely, and all three gauge
bosons acquire masses.  As a second example, consider $SU(2)$ with a scalar
field in the adjoint (triplet) representation.  In this case, the vacuum
expectation value leaves one $U(1)$ symmetry unbroken, so one gauge boson
remains massless while the other two acquire masses.  As in our general
analysis, gauge bosons corresponding to unbroken symmetry generators remain
massless.

It is interesting that this model contains both massive and massless gauge
bosons, with the distinction between these bosons created by spontaneous
symmetry breaking.  If we interpret the massive bosons as $W$ bosons and the
massless gauge boson as the photon, it is tempting to interpret this theory as
a unified model of weak and electromagnetic interactions.  Georgi and Glashow
proposed this model as a serious candidate for the theory of weak
interactions.%
\footnote{H.~Georgi and S.~L.~Glashow, \emph{Phys.\ Rev.\ Lett.} \textbf{28},
1494 (1972).}  However, Nature chooses a different model, which we will
discuss in the next section.

We turn next to a more complicated example.  Consider an $SU(3)$ gauge
theory with a scalar field in the adjoint representation.  The covariant
derivative of $\phi$ takes the form
\begin{equation}
  D_\mu\phi^a = \partial_\mu\phi^a + gf^{abc}A^b_\mu\phi^c ,
  \label{eq:20.32}
\end{equation}
and so the gauge field masses arise from the term
\begin{equation}
  \Delta\mathcal{L} = \frac{1}{2}g^2\bigl( f^{abc}A^b_\mu\phi^c \bigr)^2 .
  \label{eq:20.33}
\end{equation}
We can write this more clearly by defining the quantity
\begin{equation}
  \Phi = \phi^a t^a ,
  \label{eq:20.34}
\end{equation}
where $t^a$ are the $3\times3$ traceless Hermitian matrices that represent the
generators of $SU(3)$.  Using this notation and the definition \eqref{eq:15.68}
of the structure constants, we can rewrite the mass term \eqref{eq:20.33} as
\begin{equation}
  \Delta\mathcal{L} = -g^2\,\mathrm{tr}\Bigl[ [t^a,\Phi][t^b,\Phi] \Bigr]\,
    A^a_\mu A^{b\mu} .
  \label{eq:20.35}
\end{equation}
Now let $\Phi$ acquire a vacuum expectation value
\begin{equation}
  \langle\Phi\rangle = \Phi_0 .
  \label{eq:20.36}
\end{equation}
Since $\Phi_0$ is a traceless Hermitian matrix, we should analyze its effects
by diagonalizing it.  In principle, $\Phi_0$ could have three arbitrary
eigenvalues that sum to zero.  However, when one minimizes explicit potential
energy functions, one often finds expectation values that preserve some of the
original symmetry.  We will consider two examples.

First, $\Phi_0$ might have the orientation
\begin{equation}
  \Phi_0 = |\phi|\cdot\frac{1}{\sqrt{15}}\,
    \mathrm{diag}(1, 1, -2) .
  \label{eq:20.37}
\end{equation}
This matrix commutes with the four $SU(3)$ generators $t^1$, $t^2$, $t^3$, and
$t^8$.  Thus, the expectation value \eqref{eq:20.37} breaks $SU(3)$
spontaneously to $SU(2)\times U(1)$ and leaves the gauge bosons corresponding
to these four generators massless.  The remaining four generators of $SU(3)$
acquire the masses
\begin{equation}
  m^2 = \bigl( 3g|\phi| \bigr)^2 ,
  \label{eq:20.40}
\end{equation}
as one can check by substituting these matrices into Eq.~\eqref{eq:20.35}.

Another possible orientation for $\Phi_0$ is
\begin{equation}
  \Phi_0 = |\phi|\cdot\frac{1}{\sqrt{6}}\,
    \mathrm{diag}(1, -1, 0) .
  \label{eq:20.41}
\end{equation}
In this case, only $t^3$ and $t^8$ commute with $\Phi_0$, so the original
$SU(3)$ symmetry is broken down to $U(1)\times U(1)$.  By substituting into
\eqref{eq:20.35}, one can determine that the gauge bosons corresponding to
the remaining generators of $SU(3)$ acquire the masses
\begin{equation}
  t^1, t^2:\ m^2 = \bigl( 2g|\phi| \bigr)^2;\qquad
  t^4, t^5, t^6, t^7:\ m^2 = \bigl( g|\phi| \bigr)^2 .
  \label{eq:20.42}
\end{equation}
Still larger symmetry groups offer a wider variety of symmetry-breaking
patterns, and more complex mass matrices.  We consider one further example
in Problem~20.1.

\subsection{Formal Description of the Higgs Mechanism}

Up to this point, our study of the Higgs mechanism has been based on the
explicit analysis of scalar field Lagrangians coupled to gauge fields.  Scalar
field theories provide the simplest examples of systems with spontaneous
symmetry breaking, and the explicit calculations they allow are useful for
visualization.  But symmetries can be broken in other ways.  In the theory of
superconductivity, for example, the Abelian gauge invariance of
electromagnetism is broken by pairs of electrons that condense in the ground
state of a metal.  In Section~19.3, we argued that, in the approximation that
quark masses are very small, QCD possesses global symmetries that are
spontaneously broken by a condensation of quark-antiquark pairs.  In these
examples, spontaneous symmetry breaking is the result of strong interactions
beyond perturbation theory.  We would like to understand whether these more
general mechanisms of spontaneous symmetry breaking can also give mass to
vector bosons, and, if so, how the masses can be calculated.

To carry out this analysis, we will need to abstract several ideas from the
preceding discussion.  First, we will discuss in general terms the relations
between gauge bosons, Goldstone bosons, and global symmetry currents.  Then
we will use this information to construct the gauge boson mass matrix without
making direct use of the Lagrangian.

Consider, first, an arbitrary quantum field theory $\mathcal{L}_0$ with a
global symmetry $G$.  In Section~9.6, we derived the Noether current
corresponding to the $G$ symmetry by varying the Lagrangian by a local gauge
transformation with infinitesimal parameter $\alpha^a(x)$.  Transforming with
a constant $\alpha^a$ should leave $\mathcal{L}_0$ unchanged.  Then the more
general variation of $\mathcal{L}_0$ must take the form
\begin{equation}
  \delta\mathcal{L}_0 = \partial_\mu\alpha^a\,J^{\mu a} ,
  \label{eq:20.43}
\end{equation}
for some set of vector operators $J^{\mu a}$ built from the fields of
$\mathcal{L}_0$.  The variational principle then tells us that
\begin{equation}
  \partial_\mu J^{\mu a} = 0 .
  \label{eq:20.44}
\end{equation}
We can identify the $J^{\mu a}$ as the Noether currents of the global gauge
symmetry.

We can now couple this globally symmetric theory to non-Abelian gauge
fields, promoting the global symmetry to a local symmetry.  To first order in
$g$, the new Lagrangian should take the form
\begin{equation}
  \mathcal{L} = \mathcal{L}_0 - gA^a_\mu J^{\mu a} + \mathcal{O}(A^2) .
  \label{eq:20.45}
\end{equation}
To check this, note that the transformation \eqref{eq:20.43} compensates the
variation due to a gauge transformation of $A^a_\mu$, Eq.~\eqref{eq:15.46}, to
leading order in $g$.  The terms of order $A^2$ and higher can in general be
arranged to compensate the higher-order terms in the gauge transformation.
Thus, matrix elements involving only one insertion of the gauge field can be
evaluated using properties of the Noether currents of the original globally
symmetric theory.  Note in particular that the conservation law for these
currents, Eq.~\eqref{eq:20.44}, guarantees that the Ward identities for these
matrix elements are satisfied.

If the global symmetry of the theory $\mathcal{L}_0$ is spontaneously broken,
this theory will contain Goldstone bosons, which will stand in a special
relation to the Noether currents.  At long wavelength, the Goldstone bosons
become infinitesimal symmetry rotations of the vacuum, $Q^a|0\rangle$, where
$Q^a$ is the global charge associated with $J^{\mu a}$.  Thus, the operators
$J^{\mu a}$ have the correct quantum numbers to create Goldstone bosons from
the vacuum.  Let $|\pi^k(p)\rangle$ denote a Goldstone boson state.  In
general, there will be a current $J^{\mu a}$ that can create or destroy this
boson; we can parametrize the corresponding matrix element as
\begin{equation}
  \langle 0|J^{\mu a}(x)|\pi^k(p)\rangle = -ip^\mu F^{ak}\ e^{-ip\cdot x} ,
  \label{eq:20.46}
\end{equation}
where $p^\mu$ is the on-shell momentum of the boson and $F^{ak}$ is a matrix
of constants.  The elements $F^{ak}$ vanish when $a$ denotes a generator of
an unbroken symmetry.  Then the nonvanishing matrix elements of $F^{ak}$
connect the currents of the spontaneously broken symmetries to their
corresponding Goldstone bosons.  Since the currents $J^{\mu a}$ are conserved,
we find
\begin{equation}
  0 = \partial_\mu\langle 0|J^{\mu a}(x)|\pi^k(p)\rangle
  = -p^2 F^{ak}\,e^{-ip\cdot x} ,
  \label{eq:20.47}
\end{equation}
which implies that the bosons with nonzero matrix elements \eqref{eq:20.46}
satisfy $p^2 = 0$ on shell and so are massless.  This is another proof of
Goldstone's theorem.%
\footnote{A special case of this argument appeared in the discussion of
Eq.~\eqref{eq:19.88}.}

Since the scalar field theory that we examined earlier in this section should
be a special case of this analysis, we should find there an example of the
relation given in Eq.~\eqref{eq:20.46}.  Comparing Eqs.~\eqref{eq:20.15} and
\eqref{eq:20.45}, we see that, for the scalar field theory,
\begin{equation}
  J^{\mu a} = \partial^\mu\phi_i\,(T^a\phi)_i ,
  \label{eq:20.48}
\end{equation}
which is indeed the Noether current corresponding to the global gauge
symmetry.  Inserting the vacuum expectation value \eqref{eq:20.16}, we find
\begin{equation}
  J^{\mu a} = \partial^\mu\phi_i\,(T^a\phi_0)_i ,
  \label{eq:20.49}
\end{equation}
which leads to the set of matrix elements
\begin{equation}
  \langle 0|J^{\mu a}(x)|\pi^i(p)\rangle = -ip^\mu\,(T^a\phi_0)_i\,
    e^{-ip\cdot x} .
  \label{eq:20.50}
\end{equation}
Using this relation, we can identify
\begin{equation}
  F^{ai} = (T^a\phi_0)_i
  \label{eq:20.51}
\end{equation}
for the Higgs mechanism in a weakly coupled scalar field theory.  To be more
precise, the index $i$ runs over all components of the scalar field $\phi$.
However, we saw in the discussion below Eq.~\eqref{eq:20.19} that
\eqref{eq:20.51} is nonzero only for components $\phi_i$ that are Goldstone
bosons, and only for symmetry generators $a$ that are spontaneously broken.
Thus the nonzero components of \eqref{eq:20.51} form precisely the structure
\eqref{eq:20.46}.

As a concrete illustration of the way that the objects $T^a\phi_0$ link
spontaneously broken generators and Goldstone bosons, consider the situation
of $SU(2)$ symmetry broken by a scalar field in the vector representation.
According to the geometric picture, rotations about the 1 axis tip the vacuum
expectation value of $\phi$ into the 2 direction, rotations about the 2 axis
tip this expectation value into the 1 direction, and rotations about 3 leave
$\langle\phi\rangle$ invariant.  Thus the gauge generators $T^1$ and $T^2$ are
spontaneously broken, and the scalar field components $\phi_2$ and $\phi_1$
are the corresponding Goldstone bosons.  This accords with the result of
computing the elements of $(T^a\phi_0)_i$ explicitly: Using
$(T^a)_{bc} = \epsilon^{bac}$, we find
\begin{equation}
  (T^a\phi_0)_b = \epsilon^{bac}(\phi_0)_c = v\,\epsilon^{ba3} .
  \label{eq:20.52}
\end{equation}
Inserting this result into formula \eqref{eq:20.50}, we see that the current
of each spontaneously broken symmetry creates and destroys its own Goldstone
boson.

Now we can use this formalism to study the working of the Higgs mechanism in
this general context.  Consider the original theory $\mathcal{L}_0$ coupled to
gauge bosons of $G$.  To see how the Higgs mechanism operates, we must
compute the vacuum polarization amplitude.  This amplitude is required by the
Ward identity to be transverse, so it is necessarily of the form
\begin{equation}
  \Pi^{\mu\nu ab}(k) = i\Bigl( g^{\mu\nu} - \frac{k^\mu k^\nu}{k^2}
    \Bigr)\ \bigl( m^2_{ab} + \mathcal{O}(k^2) \bigr) .
  \label{eq:20.53}
\end{equation}
It is not easy to compute the nonsingular terms in \eqref{eq:20.53} in this
general situation, but it is straightforward to compute the singular term,
which comes from contributions with an intermediate Goldstone boson.
Combining Eqs.~\eqref{eq:20.45} and \eqref{eq:20.46}, we see that the
amplitude for a gauge boson to convert to a Goldstone boson is
\begin{equation}
  \text{($A^a$-$\pi^j$ vertex)} = -g k^\mu F^{aj} .
  \label{eq:20.54}
\end{equation}
Then the pole contribution to the vacuum polarization is
\begin{equation}
  \text{(pole contribution)} = \bigl( gk^\mu F^{aj} \bigr)\
    \frac{i}{k^2}\ \bigl( -gk^\nu F^{bj} \bigr) .
  \label{eq:20.55}
\end{equation}
Comparing \eqref{eq:20.55} with Eq.~(20.21), we identify
\begin{equation}
  m^2_{ab} = g^2\,F^{aj}F^{bj} .
  \label{eq:20.56}
\end{equation}
Notice that, in the case in which the symmetry is broken by a scalar field,
this result reverts to \eqref{eq:20.18}.  However, Eq.~\eqref{eq:20.56}
applies to any theory of spontaneously broken symmetry, whether the symmetry
breaking is apparent from the Lagrangian or whether it requires strong
interactions or other nonperturbative effects.  It is a general result, then,
that any gauge boson coupled to the current of a spontaneously broken
symmetry acquires a mass.

\section{The Glashow-Weinberg-Salam Theory of Weak Interactions}

The gauge theory of the weak and electromagnetic interactions, due to Glashow,
Weinberg, and Salam (GWS), is based on the gauge group $SU(2)\times U(1)$.  We
begin with the weak isospin $SU(2)$ and the weak hypercharge $U(1)$.  The
gauge bosons of $SU(2)$ are $W^{1\mu}$, $W^{2\mu}$, $W^{3\mu}$, with coupling
constant $g$; the gauge boson of $U(1)$ is $B^\mu$, with coupling constant
$g'$.  The Lagrangian of the gauge fields is
\begin{equation}
  \mathcal{L}_{\mathrm{gauge}} = -\frac{1}{4}W^{a\mu\nu}W^a_{\mu\nu}
    - \frac{1}{4}B^{\mu\nu}B_{\mu\nu} ,
  \label{eq:20.58}
\end{equation}
where
\begin{equation}
  W^{a\mu\nu} = \partial^\mu W^{a\nu} - \partial^\nu W^{a\mu}
    + g\epsilon^{abc}W^{b\mu}W^{c\nu},\qquad
  B^{\mu\nu} = \partial^\mu B^\nu - \partial^\nu B^\mu .
\end{equation}
The covariant derivative acting on a field with weak isospin representation
matrices $T^a$ and weak hypercharge $Y$ is
\begin{equation}
  D_\mu = \partial_\mu - igW^a_\mu T^a - ig'B_\mu Y .
  \label{eq:20.65}
\end{equation}

The symmetry is spontaneously broken by a complex scalar doublet
\begin{equation}
  \phi = \begin{pmatrix}\phi^+\\ \phi^0\end{pmatrix}
\end{equation}
with hypercharge $Y = 1/2$.  The scalar potential is
\begin{equation}
  V(\phi) = \mu^2\phi^{\dagger}\phi + \lambda(\phi^{\dagger}\phi)^2 ,
\end{equation}
with $\mu^2 < 0$, so that $\phi$ acquires a vacuum expectation value
\begin{equation}
  \langle\phi\rangle = \frac{1}{\sqrt{2}}\begin{pmatrix}0\\ v\end{pmatrix} .
  \label{eq:20.58b}
\end{equation}
Substituting this into the scalar kinetic energy term, we find the gauge
boson mass terms.  The charged gauge bosons
\begin{equation}
  W^{\pm\mu} = \frac{1}{\sqrt{2}}\bigl( W^{1\mu} \mp iW^{2\mu} \bigr)
  \label{eq:20.63}
\end{equation}
acquire the masses
\begin{equation}
  m_W = \frac{1}{2}gv .
\end{equation}
The neutral gauge bosons mix.  The photon and the $Z^0$ are the mass
eigenstates
\begin{equation}
  A_\mu = B_\mu\cos\theta_W + W^3_\mu\sin\theta_W,\qquad
  Z_\mu = -B_\mu\sin\theta_W + W^3_\mu\cos\theta_W ,
  \label{eq:20.64}
\end{equation}
where the weak mixing angle $\theta_W$ is defined by
\begin{equation}
  \sin\theta_W = \frac{g'}{\sqrt{g^2+g'^2}},\qquad
  \cos\theta_W = \frac{g}{\sqrt{g^2+g'^2}} .
\end{equation}
The photon remains massless, and the $Z^0$ acquires the mass
\begin{equation}
  m_Z = \frac{1}{2}\sqrt{g^2+g'^2}\,v = \frac{m_W}{\cos\theta_W} .
  \label{eq:20.73}
\end{equation}

We can rewrite the covariant derivative \eqref{eq:20.65} in the form
\begin{equation}
  D_\mu = \partial_\mu - \frac{ig}{\sqrt{2}}\bigl( W^{+\mu}T^+ + W^{-\mu}T^-
    \bigr) - \frac{ig}{\cos\theta_W}\ Z^\mu\bigl( T^3 - Q\sin^2\theta_W
      \bigr) - ieQA_\mu ,
  \label{eq:20.71}
\end{equation}
where
\begin{equation}
  e = g\sin\theta_W ,
  \label{eq:20.72}
\end{equation}
and $Q = T^3 + Y$ is the electric charge.  We see here that the couplings of
all of the weak bosons are described by two parameters: the well-measured
electron charge $e$, and a new parameter $\theta_W$.  The couplings induced by
$W$ and $Z$ exchange will also involve the masses of these particles.  However,
these masses are not independent, since it follows from Eqs.~\eqref{eq:20.63}
that
\begin{equation}
  m_W = m_Z\cos\theta_W .
  \label{eq:20.73b}
\end{equation}
All effects of $W$ and $Z$ exchange processes, at least at tree level, can be
written in terms of the three basic parameters $e$, $\theta_W$, and $m_W$.

\subsection{Coupling to Fermions}

The covariant derivative \eqref{eq:20.71} uniquely determines the coupling of
the $W$ and $Z^0$ fields to fermions, once the quantum numbers of the fermion
fields are specified.  To determine these quantum numbers, we must take
account of the fact, mentioned in Section~17.3, that the $W$ boson couples
only to left-handed helicity states of quarks and leptons.

At the level of the classical Lagrangian, there is no difficulty in
constructing theories in which the left- and right-handed components of a
fermion field couple differently to gauge bosons.%
\footnote{In Section 19.4, we argued that there is a possible problem with
this strategy at the level of quantum corrections.  We will check below
whether the specific model we construct avoids this problem.}  Already in
Section~3.2 we saw that the kinetic energy term for Dirac fermions splits
into separate pieces for the left- and right-handed fields:
\begin{equation}
  \bar{\psi}i\slashed{\partial}\psi = \bar{\psi}_L i\slashed{\partial}\psi_L
    + \bar{\psi}_R i\slashed{\partial}\psi_R .
  \label{eq:20.74}
\end{equation}
When we couple $\psi$ to a gauge field, we can assign $\psi_L$ and $\psi_R$ to
different representations of the gauge group.  Then the two terms on the
right-hand side of \eqref{eq:20.74} will contain two different covariant
derivatives, and these will imply two different sets of couplings.

In the GWS model, we can use this technique to insure that only the
left-handed components of the quark and lepton fields couple to the $W$
bosons.  We assign the left-handed fermion fields to doublets of $SU(2)$,
while making the right-handed fermion fields singlets under this group.
Once we have specified the $T^3$ value for each fermion field, the value of
$Y$ that we must assign follows from $Q = T^3 + Y$.  This means that the $Y$
assignments will also be different for the left- and right-handed components
of quarks and leptons.  For the right-handed fields, $T^3 = 0$, and so we
reproduce the standard electric charges by assigning $Y$ to equal the
electric charge.  For example, for the right-handed $u$ quark field,
$Y = +2/3$; for $e_R$, $Y = -1$.  For the left-handed fields,
\begin{equation}
  E_L = \begin{pmatrix}\nu_L\\ e_L\end{pmatrix},\qquad
  Q_L = \begin{pmatrix}u_L\\ d_L\end{pmatrix} ,
  \label{eq:20.75}
\end{equation}
the assignments $Y = -1/2$ and $Y = +1/6$, respectively, combine with
$T^3 = \pm 1/2$ to give the correct electric charge assignments.  Since the
left- and right-handed fermions live in different representations of the
fundamental gauge group, it is often useful to think of these components as
distinct particles, which are mixed by the fermion mass terms.

In fact, the construction of fermion mass terms is a serious problem, because
all possible such terms are forbidden by global gauge invariances.  For
example, we cannot write an electron mass term
\begin{equation}
  \Delta\mathcal{L} = -m_e\bigl( \bar{e}_Le_R + \bar{e}_Re_L \bigr) ,
  \label{eq:20.76}
\end{equation}
because the fields $e_L$ and $e_R$ belong to different $SU(2)$
representations and have different $U(1)$ charges.  For the next few pages, we
will ignore this problem by treating all fermion fields as massless.  This
description will suffice to analyze the structure of the weak interactions at
high energies, where the quark and lepton masses can be ignored.  At the end
of this section we will return to the problem of writing quark and lepton
mass terms in the GWS theory.  The solution to this problem will reinforce the
idea that the left- and right-handed fermion fields are fundamentally
independent entities, mixed to form massive fermions by some subsidiary
process.

If we ignore fermion masses, the Lagrangian for the weak interactions of
quarks and leptons follows directly from the charge assignments given above.
The fermion kinetic energy terms for $e$, $\nu$, $u$, and $d$ are
\begin{equation}
  \mathcal{L} = \bar{E}_L(i\slashed{D})E_L + \bar{e}_R(i\slashed{D})e_R
    + \bar{Q}_L(i\slashed{D})Q_L + \bar{u}_R(i\slashed{D})u_R
    + \bar{d}_R(i\slashed{D})d_R .
  \label{eq:20.77}
\end{equation}
In each term, the covariant derivative is given by Eq.~\eqref{eq:20.65}, with
$T^a$ and $Y$ evaluated in the particular representation to which that
fermion field belongs.  For example,
\begin{equation}
  \bar{Q}_L(i\slashed{D})Q_L = \bar{Q}_L i\gamma^\mu\Bigl( \partial_\mu
    - igW^a_\mu T^a - i\frac{g'}{6}B_\mu \Bigr)Q_L .
  \label{eq:20.78}
\end{equation}
A right-handed neutrino would have zero coupling both to $SU(2)$ and to
$U(1)$, so we have simply omitted this field from Eq.~\eqref{eq:20.77}.

To work out the physical consequences of the fermion-vector boson couplings,
we should write Eq.~\eqref{eq:20.77} in terms of the vector boson mass
eigenstates, using the form of the covariant derivative given in
Eq.~\eqref{eq:20.71}.  Equation \eqref{eq:20.77} then takes the form
\begin{equation}
  \begin{aligned}
    \mathcal{L} &= \bar{E}_L(i\slashed{\partial})E_L + \bar{e}_R(i\slashed{\partial})e_R
      + \bar{Q}_L(i\slashed{\partial})Q_L + \bar{u}_R(i\slashed{\partial})u_R
      + \bar{d}_R(i\slashed{\partial})d_R\\
    &\quad + g\bigl( W^+_\mu J^{\mu+} + W^-_\mu J^{\mu-} + Z^0_\mu J^\mu_Z
      \bigr) + eA_\mu J^\mu_{EM} ,
  \end{aligned}
  \label{eq:20.79}
\end{equation}
where
\begin{equation}
  \begin{aligned}
    J^{\mu+}_W &= \frac{1}{\sqrt{2}}\bigl( \bar{\nu}_L\gamma^\mu e_L
      + \bar{u}_L\gamma^\mu d_L \bigr),\\
    J^{\mu-}_W &= \frac{1}{\sqrt{2}}\bigl( \bar{e}_L\gamma^\mu\nu_L
      + \bar{d}_L\gamma^\mu u_L \bigr),\\
    J^\mu_Z &= \frac{1}{\cos\theta_W}\ \Bigl[ \bar{\nu}_L\gamma^\mu
      \Bigl( \frac{1}{2} \Bigr)\nu_L + \bar{e}_L\gamma^\mu\Bigl( -\frac{1}{2}
        + \sin^2\theta_W \Bigr)e_L + \bar{e}_R\gamma^\mu(\sin^2\theta_W)e_R\\
    &\quad + \bar{u}_L\gamma^\mu\Bigl( \frac{1}{2} - \frac{2}{3}\sin^2\theta_W
      \Bigr)u_L + \bar{u}_R\gamma^\mu\Bigl( -\frac{2}{3}\sin^2\theta_W
        \Bigr)u_R\\
    &\quad + \bar{d}_L\gamma^\mu\Bigl( -\frac{1}{2} + \frac{1}{3}\sin^2\theta_W
      \Bigr)d_L + \bar{d}_R\gamma^\mu\Bigl( \frac{1}{3}\sin^2\theta_W
        \Bigr)d_R \Bigr],\\
    J^\mu_{EM} &= \bar{e}\gamma^\mu(-1)e + \bar{u}\gamma^\mu\Bigl(
      +\frac{2}{3} \Bigr)u + \bar{d}\gamma^\mu\Bigl( -\frac{1}{3} \Bigr)d .
  \end{aligned}
  \label{eq:20.80}
\end{equation}
Here we have used Eq.~(20.67) to simplify the $W$ boson currents.
Notice that the current $J^\mu_{EM}$ associated with the photon field is
indeed the standard electromagnetic current.

\subsection{Anomaly Cancellation}

As we have just seen, there is no difficulty in writing a Lagrangian that
couples the GWS gauge bosons to fermions in a chiral fashion.  However, these
chiral couplings do present a potential problem that appears at the level of
one-loop corrections.  In Section~19.2, we saw that an axial current that is
conserved at the level of the classical equations of motion can acquire a
nonzero divergence through one-loop diagrams that couple this current to a
pair of gauge bosons.  The Feynman diagram that contains this anomalous
contribution is a triangle diagram with the axial current and the two gauge
currents at its vertices.  In a gauge theory in which gauge bosons couple to a
chiral current, the dangerous triangle diagrams appear in the one-loop
corrections to the three-gauge-boson vertex function.  The anomalous terms
violate the Ward identity for this amplitude.  Thus, as we argued in
Section~19.4, theories in which gauge bosons couple to chiral currents can be
gauge invariant only if the anomalous contribution somehow disappears.
Fortunately, as we saw there, the anomalous terms can be arranged to cancel
when one sums over all possible fermion species that can circulate in these
diagrams.%
\footnote{If you have not read Chapter 19, but you are willing to assume that
the fermion triangle diagram contains a contribution that violates gauge
invariance, you should still be able to follow the argument that follows.}

Within the GWS theory, the requirement from experiment that the weak
interaction currents are left-handed forced us to choose a chiral gauge
coupling.  Now we must check that the anomalous terms from the triangle
diagrams cancel as required.  We will find that they do, but only through a
subtle and rather magical interplay of the quantum numbers of quarks and
leptons.

The anomalous term of a triangle diagram of three gauge bosons $A^a$, $A^b$,
and $A^c$ is proportional to the group theoretic invariant
\begin{equation}
  \mathrm{tr}\Bigl[ \gamma^5\ t^a\{t^b, t^c\} \Bigr] ,
  \label{eq:20.81}
\end{equation}
where the trace is taken over all fermion species.  The anticommutator comes
from taking the sum of two triangle diagrams in which the fermions circle in
opposite directions.  The factor $\gamma^5$ registers the fact that the
anomaly is associated with chiral currents; this factor equals $-1$ for
left-handed fermions and $+1$ for right-handed fermions.  In theories such as
QED or QCD in which the gauge bosons couple equally to right- and left-handed
species, the anomalies automatically cancel.  This bookkeeping method is a
special case of the more general method presented in Section~19.4.

To evaluate the anomalies of the GWS theory, it is easiest to work in the
basis of $SU(2)\times U(1)$ gauge bosons, before the mixing to the photon and
$Z^0$ mass eigenstates.  It suffices to evaluate the triangle diagrams for
massless fermions, so that right- and left-handed fermions have distinct
quantum numbers.  However, we must consider not only the anomalies of
diagrams with three $SU(2)\times U(1)$ gauge bosons, but also diagrams with
both weak-interaction gauge bosons and color $SU(3)$ gauge bosons of QCD.
If we consider effects of gravity on the weak-interaction gauge theory, there
is also a possibly anomalous diagram with one weak-interaction gauge boson
and two gravitons.  We can omit diagrams, such as the anomaly of three
$SU(3)$ bosons or of one $SU(3)$ boson and two gravitons, in which all of the
couplings are left-right symmetric.  Then the full set of diagrams with
possible anomalous terms is shown in Fig.~\ref{fig:20.2}.  All of the possible
anomalies must cancel if the Ward identities of the $SU(2)\times U(1)$ gauge
theory are to be satisfied.

\begin{figure}[t]
  \centering
  \includegraphics[width=0.6\textwidth]{fig20.2.png}
  \caption{Triangle diagrams that can potentially generate anomalies in the
    GWS theory.}
  \label{fig:20.2}
\end{figure}

It is a special property of $SU(2)$ gauge theory that the anomaly of three
$SU(2)$ gauge bosons always vanishes; this result follows from the property
of Pauli sigma matrices $\{\sigma^a,\sigma^b\} = 2\delta^{ab}$, which implies
that the trace \eqref{eq:20.81} vanishes.  The anomalies containing one
$SU(3)$ boson or one $SU(2)$ boson are proportional to
\begin{equation}
  \mathrm{tr}[t^a] = 0
  \qquad\text{or}\qquad
  \mathrm{tr}[\tau^a] = 0 .
  \label{eq:20.82}
\end{equation}
The remaining nontrivial anomalies are those of one $U(1)$ boson with two
$SU(2)$ bosons or two $SU(3)$ bosons, the anomaly of three $U(1)$ bosons, and
the gravitational anomaly with one $U(1)$ gauge boson.

The anomaly of one $U(1)$ boson with two $SU(3)$ bosons is proportional to
the group theory factor
\begin{equation}
  \mathrm{tr}\bigl[ t^at^bY \bigr] = \frac{1}{2}\delta^{ab}\cdot
    \sum_{\text{quarks}} Y ,
  \label{eq:20.83}
\end{equation}
where the sum runs over left-handed quarks and right-handed quarks, with an
extra $(-1)$ for the left-handed contributions.  Inserting the charge
assignments given above for $u_L$, $d_L$, $u_R$, and $d_R$, we find
\begin{equation}
  \sum_{\text{quarks}} Y = -2\cdot\frac{1}{6} + \Bigl( \frac{2}{3} \Bigr)
    + \Bigl( -\frac{1}{3} \Bigr) = 0 .
  \label{eq:20.84}
\end{equation}
The anomaly of one $U(1)$ boson with two $SU(2)$ bosons is proportional to
the sum of hypercharges of the left-handed doublets.  Adding the lepton
doublet and the quark doublet, we find
\begin{equation}
  \sum_{\text{doublets}} Y = 3\cdot\Bigl( \frac{1}{6} \Bigr)
    + \Bigl( -\frac{1}{2} \Bigr) = 0 .
\end{equation}
The anomaly of three $U(1)$ bosons is proportional to the sum of the cubes of
the hypercharges.  Adding the contributions of all left- and right-handed
fermions with the appropriate signs, one finds that this also vanishes, with
the color factor of 3 for the quarks:
\begin{equation}
  \sum_{\text{all fermions}} Y^3 = 3\cdot\Bigl[ -2\Bigl( \frac{1}{6} \Bigr)^3
    + \Bigl( \frac{2}{3} \Bigr)^3 + \Bigl( -\frac{1}{3} \Bigr)^3 \Bigr]
    - 2\Bigl( -\frac{1}{2} \Bigr)^3 + (-1)^3 = 0 .
  \label{eq:20.85}
\end{equation}
Finally, the gravitational anomaly with one $U(1)$ gauge boson is
proportional to the sum of the hypercharges, which we have already shown to
vanish.  Thus, the GWS theory is anomaly free.

The cancellation of anomalies in the GWS theory is a remarkable and
nontrivial consistency condition.  It required the quarks to have fractional
charges, the leptons to have integer charges, and the number of colors to be
3.  The fact that these conditions are all satisfied by the observed spectrum
of particles is a strong indication that the GWS theory is on the right
track.

\subsection{Experimental Consequences of the Glashow-Weinberg-Salam Theory}

The GWS theory makes a number of precise predictions for weak interaction
phenomena.  At low energies, the exchange of the heavy $W$ and $Z$ bosons
gives rise to effective four-fermion interactions.  The charged-current
interactions reproduce the familiar structure of the Fermi theory of weak
interactions, with the Fermi constant given by
\begin{equation}
  \frac{G_F}{\sqrt{2}} = \frac{g^2}{8m_W^2} .
  \label{eq:20.91}
\end{equation}
The neutral-current interactions give new predictions that are not present in
the Fermi theory.  These predictions have been tested in a wide variety of
experiments.

Further tests of the GWS theory are available at higher energies.  The
process $e^{+}e^{-}\to f\bar{f}$ is affected in an essential way, since the
theory contains a new diagram with $s$-channel $Z^0$ exchange, which
interferes with the standard photon exchange diagram, as shown in
Fig.~\ref{fig:20.4}.  It is straightforward to work out the effects of this
interference using the methods of Section~5.2, so we have left this analysis
as Problem~20.3.

\begin{figure}[t]
  \centering
  \includegraphics[width=0.55\textwidth]{fig20.4.png}
  \caption{Diagrams contributing to the process $e^{+}e^{-}\to f\bar{f}$ in
    the Glashow-Weinberg-Salam theory.}
  \label{fig:20.4}
\end{figure}

As the center-of-mass energy approaches $m_Z$, the $Z^0$ appears directly as
a resonance in the $e^{+}e^{-}$ annihilation cross section.  Similarly, both
the $W$ and the $Z$ can be observed as resonances in quark-antiquark
annihilation, viewed as a parton subprocess in proton-antiproton scattering.
The positions of these resonances are predicted from $G_F$, $\sin^2\theta_W$,
and the value of $e$ or $\alpha$, according to Eqs.~\eqref{eq:20.72} and
\eqref{eq:20.91}.  Using these relations, we find
\begin{equation}
  m_W = \sqrt{\frac{\alpha\pi}{\sqrt{2}G_F\sin^2\theta_W}},\qquad
  m_Z = \sqrt{\frac{\alpha\pi}{\sqrt{2}G_F\sin^2\theta_W\cos^2\theta_W}} .
  \label{eq:20.95}
\end{equation}
The detailed shape of the $Z^0$ resonance is shown in Fig.~\ref{fig:20.5}.
The experimental measurements shown are compared to a theoretical curve with
the resonance position adjusted for the best fit.  The height and width of
the resonance are then predicted by the GWS theory.  The resonance is
broadened to higher energies by processes in which the electron and positron
radiate collinear photons before annihilation; this correction was discussed
in Problem~5.5.

\begin{figure}[t]
  \centering
  \includegraphics[width=0.6\textwidth]{fig20.5.png}
  \caption{The total cross section for $e^{+}e^{-}$ annihilation to hadrons
    for $E_{cm}$ close to the $Z^0$ boson mass, as measured by the ALEPH,
    DELPHI, L3, and OPAL experiments and compiled by the Particle Data Group,
    \emph{Phys.\ Rev.\ D} \textbf{50}, (1994), Fig.~32.14.  References to the
    original articles are given there.  The solid curve is the prediction of
    the GWS theory.}
  \label{fig:20.5}
\end{figure}

Because the Lagrangian of the GWS theory treats left- and right-handed
fermions as distinct species with completely different quantum numbers, the
couplings of the $Z^0$ to left- and right-handed fermions differ
significantly.  One manifestation of this is the presence of a polarization
asymmetry, a net polarization of fermions produced in the decay
$Z^0\to f\bar{f}$, or an asymmetry in the inverse process of $Z^0$
production.  This asymmetry can be read directly from the form of the $Z^0$
current given in \eqref{eq:20.80}:
\begin{equation}
  \mathcal{A}_{LR} = \frac{\Gamma(Z^0\to f_L\bar{f}_R)
    - \Gamma(Z^0\to f_R\bar{f}_L)}
    {\Gamma(Z^0\to f_L\bar{f}_R) + \Gamma(Z^0\to f_R\bar{f}_L)}
  = \frac{\bigl( \tfrac{1}{2} - |Q_f|\sin^2\theta_W \bigr)^2
    - \bigl( Q_f\sin^2\theta_W \bigr)^2}
    {\bigl( \tfrac{1}{2} - |Q_f|\sin^2\theta_W \bigr)^2
    + \bigl( Q_f\sin^2\theta_W \bigr)^2} .
  \label{eq:20.96}
\end{equation}
For a realistic value $\sin^2\theta_W = 0.23$, this expression gives a 15\%
asymmetry for charged leptons and a 95\% asymmetry for $d$, $s$, and $b$
quarks.  The asymmetry can be checked experimentally for leptons by
measuring the polarization of $\tau$ leptons at the $Z^0$ resonance, or by
measuring the relative cross sections for producing the resonance using left-
versus right-handed electrons.  For quarks, the asymmetry can be determined
indirectly from the forward-backward production asymmetry on the resonance,
as explained in Problem~20.3.

Because the weak neutral current has so many different manifestations, the
GWS theory of weak interactions can be subjected to a stringent test by
comparing the values of the parameter $\sin^2\theta_W$ needed to account for
each of its predicted effects.  Table~\ref{tab:20.1} presents the values of
$\sin^2\theta_W$ extracted from a wide variety of weak interaction neutral
current effects and asymmetries.  In all cases, one-loop radiative corrections
must be included to analyze the experiment at the required level of accuracy.
These radiative corrections involve some subtlety.  First, one must adopt a
specific renormalization convention that defines $\sin^2\theta_W$ and use it
consistently in all calculations.  The table shows results for two different
choices of this convention.  In both conventions, the values of
weak-interaction observables are taken to be functions of $\alpha$, $G_F$,
and a third independent parameter.

\begin{table}[t]
  \centering
  \caption{Values of $\sin^2\theta_W$ from Weak Interaction Experiments}
  \begin{tabular}{lll}
    \toprule
    $\sin^2\theta^{MS}_W$ & $\sin^2\theta^{MSbar}_W$ & Observed Quantity or
      Process\\
    \midrule
    0.2247 (21) & 0.2320 (6) & $m_Z$\\
    0.2264 (25) & 0.2338 (22) & $m_W$\\
    0.2250 (18) & 0.2322 (6) & $\Gamma_Z$\\
    0.2243 (17) & 0.2315 (11) & Lepton $f$-b asymmetries at the $Z^0$\\
    0.2245 (17) & 0.2317 (8) & All pair-production asymmetries at the $Z^0$\\
    0.2221 (17) & 0.2292 (10) & $A_{LR}$ at the $Z^0$\\
    0.2260 (48) & 0.233 (5) & Deep inelastic neutrino scattering\\
    0.205 (31) & 0.212 (32) & Neutrino-proton elastic scattering\\
    0.224 (9) & 0.231 (9) & Neutrino-electron elastic scattering\\
    0.216 (8) & 0.223 (8) & Atomic parity violation\\
    0.216 (17) & 0.223 (18) & Parity violation in inelastic $e^{-}$
      scattering\\
    \bottomrule
  \end{tabular}
  \label{tab:20.1}
\end{table}

The values listed here are obtained by fitting experimental observations by
adjusting the value of $\sin^2\theta^{MS}_W$ or $\sin^2\theta^{\overline{MS}}_W$,
taking $\alpha$ and $G_F$ as accurately known parameters.  The numbers in
parentheses are the standard errors in the last displayed digits.  The
conversion from the experimentally measured quantities to
$\sin^2\theta_W$ or $\sin^2\theta^{\overline{MS}}_W$ depends on the value of
the top quark mass and the mass of the Higgs boson.  These values assume a
top quark mass of 169 GeV and a Higgs mass of 300 GeV; the quoted errors
include an uncertainty of 17 GeV in the top quark mass and a range from
60 GeV to 1000 GeV for the Higgs mass.  The differences in the relative
errors between the two columns reflect the importance of this theoretical
uncertainty.  Some observables depend weakly on $\alpha_s$; these values
assume $\alpha_s(m_Z) = 0.120 \pm 0.007$.  This table is taken from the
article of P.~Langacker and J.~Erler for the Particle Data Group,
\emph{Phys.\ Rev.\ D} \textbf{50}, 1304 (1994).  That article contains a full
set of references and a discussion of the sources of uncertainty in these
determinations.

In the first column this parameter is the mass ratio $m_W/m_Z$, and, following
the tree-level expression \eqref{eq:20.73}, we consider this ratio to define a
renormalized value of $\sin^2\theta_W$:
\begin{equation}
  \sin^2\theta^s_W = 1 - \frac{m_W^2}{m_Z^2} .
  \label{eq:20.97}
\end{equation}
In the second column, the third parameter is $\sin^2\theta_W$ computed from
the weak interaction coupling constants defined by minimal subtraction
(Eq.~\eqref{eq:11.77}).  The differences between different definitions of
$\sin^2\theta_W$ appear at the level of one-loop computations and can reveal
interesting physics; this subject is discussed in Section~21.3.

A second subtlety is that the one-loop corrections to weak neutral current
processes depend on the value of the $t$ quark mass, which has only recently
been determined and is still somewhat poorly known.  The dependence on the
$t$ quark mass is relatively strong, for interesting reasons that we will
discuss in Section~21.3.  The one-loop corrections also depend weakly on
properties of the particles responsible for the spontaneous symmetry
breaking.

We can see from Table~\ref{tab:20.1} that a wide variety of effects due to
the weak neutral current have been observed, with magnitudes accounted for by
a single, consistent value of $\sin^2\theta_W$.  This remarkable concordance
of theory and experiment gives us confidence that the Glashow-Weinberg-Salam
theory is indeed the correct description of weak and electromagnetic
interactions.

\subsection{Fermion Mass Terms}

We now return to the problem of writing mass terms for the quarks and
leptons.  Recall that one cannot put ordinary mass terms into the Lagrangian,
because the left- and right-handed components of the various fermion fields
have different gauge quantum numbers and so simple mass terms violate gauge
invariance.  To give masses to the quarks and leptons, we must again invoke
the mechanism of spontaneous symmetry breaking.

We began this section by assuming that a scalar field $\phi$ acquires a
vacuum expectation value \eqref{eq:20.58}, in order to give mass to the $W$
and $Z$ bosons.  This scalar field needed to be a spinor under $SU(2)$ and to
have $Y = 1/2$ in order to produce the correct pattern of gauge boson masses.
With these quantum numbers, we can also write a gauge-invariant coupling
linking $e_L$, $e_R$, and $\phi$, as follows:
\begin{equation}
  \Delta\mathcal{L}_e = -\lambda_e\ \bar{E}_L\cdot\phi\,e_R + \mathrm{h.c.}
  \label{eq:20.98}
\end{equation}
Here the $SU(2)$ indices of the doublets $E_L$ and $\phi$ are contracted;
notice also that the charges $Y$ of the various fields sum to zero.  The
parameter $\lambda_e$ is a new dimensionless coupling constant.  If we replace
$\phi$ in this expression by its vacuum expectation value \eqref{eq:20.58}, we
obtain
\begin{equation}
  \Delta\mathcal{L}_e = -\frac{1}{\sqrt{2}}\lambda_e\,v\ \bar{e}_Le_R
    + \mathrm{h.c.} + \cdots .
  \label{eq:20.99}
\end{equation}
This is a mass term for the electron.  The size of the mass is set by the
vacuum expectation value of $\phi$, rescaled by the new dimensionless
coupling:
\begin{equation}
  m_e = \frac{1}{\sqrt{2}}\lambda_e v .
  \label{eq:20.100}
\end{equation}
Since the electron mass is proportional to $v$, one might expect that the
masses of the electron and the $W$ boson should be of the same order.  In
fact, taking the observed values, $m_e/m_W \sim 6\times10^{-6}$.  Since
$\lambda_e$ is a renormalizable coupling, it must be treated as an input to
the theory.  Thus the GWS theory allows the electron to be very light, but it
cannot explain why the electron is so light compared to the $W$ boson.

We can write mass terms for the quark fields in the same way.  Notice that,
in the following expression, both terms are invariant under $SU(2)$ and have
zero net $Y$:
\begin{equation}
  \Delta\mathcal{L}_q = -\lambda_d\ \bar{Q}_L\cdot\phi\,d_R
    - \lambda_u\ \epsilon_{ab}\ \bar{Q}^a_L\phi^{b*}\,u_R + \mathrm{h.c.}
  \label{eq:20.101}
\end{equation}
Substituting the vacuum expectation value of $\phi$ from Eq.~\eqref{eq:20.58},
these terms become
\begin{equation}
  \Delta\mathcal{L}_q = -\frac{1}{\sqrt{2}}\lambda_d\,v\ \bar{d}_Ld_R
    - \frac{1}{\sqrt{2}}\lambda_u\,v\ \bar{u}_Lu_R + \mathrm{h.c.} + \cdots ,
  \label{eq:20.102}
\end{equation}
standard mass terms for the $d$ and $u$ quarks.  The GWS theory thus gives the
relations
\begin{equation}
  m_d = \frac{1}{\sqrt{2}}\lambda_d v,\qquad
  m_u = \frac{1}{\sqrt{2}}\lambda_u v .
  \label{eq:20.103}
\end{equation}
As with the electron, the theory parametrizes but does not explain the small
values of the $d$ and $u$ quark masses observed experimentally.

When additional generations of quarks are introduced into the theory, there
can be additional coupling terms that mix generations.  Alternatively, we can
diagonalize the Higgs couplings by choosing a new basis for the quark fields.
We will show that this is always possible in Section~20.3.  However, this
simplification of the Higgs couplings causes a complication in the gauge
couplings.  The charged-current weak interaction then couples the quark
generations through a unitary matrix, the Cabibbo-Kobayashi-Maskawa (CKM)
matrix, which is not diagonal in general.

The phenomenological properties of the Higgs boson are worked out in more
detail in the Final Project of Part~III.

\subsection{A Higgs Sector?}

Since there is no experimental evidence for the existence of the simple Higgs
boson contained in the GWS model, it is worth asking whether the $W$ and $Z$
bosons might acquire mass by a more complicated mechanism.  There are two
aspects to this question.

First, is it certain that the $W$ and $Z$ bosons are gauge bosons of a
spontaneously broken $SU(2)\times U(1)$ symmetry?  The evidence for this idea
comes from the universality of the couplings of the various quarks and leptons
to the weak interactions.  This universality is tested in the fact that the
same value of the Fermi constant describes all charged-current
weak-interaction processes, and that this same strength of coupling combined
with a single value of $\sin^2\theta_W$ describes the whole range of weak
neutral current phenomena.  We have seen, especially in the discussion of
Chapter~16, that the principle of local gauge invariance leads naturally to
the prediction of universal, flavor-independent, coupling constants.  No
other principle is known that would explain this striking regularity.  Thus
there is compelling evidence that the underlying theory of the weak
interactions is a spontaneously broken gauge theory.

However, it is quite possible that the mechanism of the spontaneous breaking
of $SU(2)\times U(1)$ is more complicated than the simple model of a single
scalar field that we have written in Eq.~(20.111).  In principle,
the breaking of $SU(2)\times U(1)$ might be the result of the dynamics of a
complicated new set of particles and interactions, which we will refer to as
the \emph{Higgs sector}.  Experiment gives us only three properties of this
new sector: First, it must generate the masses of the quarks and leptons.
Second, it must generate the masses of the $W$ and $Z$ bosons.  The third
piece of information, which is the only nontrivial one, comes from the
relation \eqref{eq:20.73} between weak boson masses in the GWS theory,
\begin{equation}
  m_W = m_Z\cos\theta_W .
  \label{eq:20.117}
\end{equation}
This relation is satisfied experimentally to better than 1\% accuracy, that
is, to the level of one-loop radiative corrections.  Whatever complicated
mechanism we invoke to generate the spontaneous breaking of $SU(2)\times U(1)$,
it should reproduce this relation in a natural way.

To understand the implications of relation \eqref{eq:20.117}, we must analyze
the gauge boson mass matrix without assuming that $SU(2)\times U(1)$ is broken
by the expectation value of a scalar field.  Actually, it is possible to
compute the gauge boson mass matrix under much less restrictive assumptions,
using the argument given at the very end of Section~20.1.  There we
constructed the gauge boson mass matrix from the matrix elements for gauge
currents to create or destroy Goldstone bosons.  We will now show that
relation \eqref{eq:20.117} follows for a large class of models of
$SU(2)\times U(1)$ breaking for which these matrix elements satisfy certain
simple properties.

Any model of weak-interaction symmetry breaking must contain some set of
fields that is responsible for the spontaneous breaking of $SU(2)\times U(1)$.
Think of this sector of the theory as a field theory with a global
$SU(2)\times U(1)$ symmetry, which is promoted to a local symmetry through its
coupling to gauge bosons.  In the theory with global symmetry, this symmetry
is spontaneously broken to $U(1)$.  Since three continuous symmetries are
spontaneously broken, this sector must supply three Goldstone bosons, which
will eventually be eaten by $W^+$, $W^-$, and $Z^0$.  Call these three bosons
$\pi^a$, where $a = 1, 2, 3$.  Let $J^{\mu A}$ be the $SU(2)$ symmetry
currents of the new sector, and let $J^{\mu Y}$ be the $U(1)$ current.  The
gauge boson mass matrix will then be constructed from the matrix elements
\eqref{eq:20.46}, which here take the form
\begin{equation}
  \langle 0|J^{\mu A}|\pi^b(p)\rangle = -ip^\mu F^{Ab} ,
  \label{eq:20.118}
\end{equation}
with $A = 1, 2, 3, Y$ and $b = 1, 2, 3$.  Using the method of
Eq.~\eqref{eq:20.55}, we find that the gauge boson vacuum polarization
contains the pole term
\begin{equation}
  -\frac{i}{k^2}\ \bigl( g_AF^{Ac} \bigr)\bigl( g_BF^{Bc} \bigr) ,
  \label{eq:20.119}
\end{equation}
summed over $c$, where $g_A = g$ for $A = 1, 2, 3$ and $g_A = g'$ for
$A = Y$.  Then we can identify the gauge boson mass matrix as
\begin{equation}
  m^2_{AB} = g_Ag_B\,F^{Ac}F^{Bc} .
  \label{eq:20.120}
\end{equation}
To reproduce the known form of the weak gauge boson mass matrix, we must now
place constraints on the $F^{Ab}$.  First we must insure that the photon
remains massless.  This follows if the linear combination of charges
Eq.~(20.69) annihilates the vacuum.  In the language of
Eq.~\eqref{eq:20.118}, we must insist that the corresponding linear
combination of currents cannot excite a Goldstone boson:
\begin{equation}
  \langle 0|\bigl( J^{\mu 3} + J^{\mu Y} \bigr)|\pi^0(p)\rangle = 0 .
  \label{eq:20.121}
\end{equation}
We can also achieve relation \eqref{eq:20.117}, using the following
additional assumption: The symmetry-breaking sector has an $SU(2)$ global
symmetry, under which the three Goldstone bosons and the three $SU(2)$ gauge
currents transform as triplets, which remains exact when the $SU(2)$ gauge
symmetry is spontaneously broken.  This global $SU(2)$ symmetry implies that,
if $A = a = 1, 2, 3$ in Eq.~\eqref{eq:20.118},
\begin{equation}
  \langle 0|J^{\mu a}|\pi^b(p)\rangle = -iF\,p^\mu\,\delta^{ab} ,
  \label{eq:20.122}
\end{equation}
where $F$ is a parameter with the dimensions of mass.  Combining
\eqref{eq:20.122} and \eqref{eq:20.121}, we have
\begin{equation}
  \langle 0|J^{\mu Y}|\pi^3(p)\rangle = +iF\,p^\mu .
  \label{eq:20.123}
\end{equation}
Inserting this form for $F^{Ab}$ into \eqref{eq:20.120}, we find the gauge
boson mass matrix
\begin{equation}
  m^2 = F^2\ \begin{pmatrix}
    g^2 & & & \\
    & g^2 & & \\
    & & g^2 & -gg'\\
    & & -gg' & g'^2
  \end{pmatrix},
  \label{eq:20.124}
\end{equation}
where the matrix acts on the gauge boson $(A^1_\mu, A^2_\mu, A^3_\mu, B_\mu)$.
The eigenvectors of this matrix are precisely \eqref{eq:20.63} and
\eqref{eq:20.64}.  To reproduce the eigenvalues, we need only define
\begin{equation}
  v = 2F .
  \label{eq:20.125}
\end{equation}
We have now shown that the GWS relation between the $W$ and $Z$ boson masses
is not special to the situation in which the gauge symmetry is broken by a
single scalar field.  This relation follows from the much more general
assumption of an unbroken global $SU(2)$ symmetry of the Higgs sector.  This
symmetry is often called \emph{custodial $SU(2)$ symmetry}.%
\footnote{P.~Sikivie, L.~Susskind, M.~Voloshin, and V.~Zakharov,
\emph{Nucl.\ Phys.\ B} \textbf{173}, 189 (1980).}  We have seen this symmetry
already as the global $SU(2)$ symmetry of the weak-interaction effective
Lagrangian Eq.~(20.94).

For the case of a single scalar field, the custodial symmetry arises in the
following way: If we write the field $\phi$ in terms of its four real
components, the Lagrangian Eq.~(20.111) (ignoring the gauge couplings)
has $O(4)$ global symmetry.  The vacuum expectation value of $\phi$ breaks
this symmetry down to $O(3)$, that is, $SU(2)$.

However, there are many other quantum field theories that break $SU(2)$
spontaneously while leaving another global $SU(2)$ symmetry unbroken.  One
rather complex example is given by QCD with two massless flavors, if we
identify the gauged $SU(2)$ with the symmetry generated by $U_L$ in
\eqref{eq:19.82} and identify the custodial $SU(2)$ with vectorial isospin
symmetry.  A copy of the familiar strong interactions with a mass scale large
enough to give $F = 125$ GeV would be a perfectly acceptable model for the
Higgs sector.  (Unfortunately, it is not easy in this model to generate masses
for the quarks and leptons.)

The question of the nature of the Higgs sector and the explicit mechanism of
$SU(2)\times U(1)$ breaking is probably the most pressing open problem in the
theory of elementary particles.  We will discuss this question further in the
Epilogue.

\section{Symmetries of the Theory of Quarks and Leptons}

Putting together the theory of strong interactions described in Chapter~17 and
the theory of weak and electromagnetic interactions described in the previous
section, we have now constructed a complete description of elementary particle
interactions.  It is interesting to investigate the symmetries of this theory,
to ask what might be the fundamental symmetries of the underlying description
of Nature.

We have already seen, in the arguments leading up to Eq.~\eqref{eq:15.17},
that the Lagrangian of a gauge theory is highly restricted by the conditions
of renormalizability and gauge invariance.  In this section, we will construct
the most general renormalizable Lagrangian consistent with the
$SU(3)\times SU(2)\times U(1)$ gauge symmetries of the strong, weak and
electromagnetic interactions.  We can then ask what further global symmetries
we must impose on this theory in order to give it the global symmetries that
we see in Nature.

As a first step, we will ignore the Higgs scalar field and the mass terms of
quarks, leptons, and gauge bosons.  Then the Lagrangian of the theory of
quarks and leptons is entirely specified by gauge invariance and
renormalizability.  We have
\begin{equation}
  \mathcal{L}_{G} = -\frac{1}{4}\sum_i F^{i\mu\nu}F^i_{\mu\nu}
    + \sum_J \bar{\psi}_J i\slashed{D}\psi_J ,
  \label{eq:20.126}
\end{equation}
where the index $i$ runs over the three factors of the gauge group and the
index $J$ runs over the various multiplets of chiral fermions.

In principle, we could add to \eqref{eq:20.126} the following pseudoscalar
pure gauge operators:
\begin{equation}
  \Delta\mathcal{L}_{\theta} = \sum_i \theta_i\ \frac{g_i^2}{32\pi^2}\,
    \epsilon^{\mu\nu\rho\sigma}F^i_{\mu\nu}F^i_{\rho\sigma} .
  \label{eq:20.127}
\end{equation}
These terms are apparently odd under both $P$ and $T$.  However, we saw at
the end of Section~19.2 that terms of this form can be generated or canceled
by making a change of variables in the effective action.  For example, the
change of variables on the right-handed electron field
\begin{equation}
  e_R \to e^{i\alpha}e_R
  \label{eq:20.128}
\end{equation}
produces, according to \eqref{eq:19.78} or \eqref{eq:19.79}, a correction to
the Lagrangian involving the $P$- and $T$-odd combination of field strengths
for the $U(1)$ gauge field
\begin{equation}
  \Delta\mathcal{L} = \alpha\,\frac{g'^2}{32\pi^2}\,
    \epsilon^{\mu\nu\rho\sigma}B_{\mu\nu}B_{\rho\sigma} .
  \label{eq:20.129}
\end{equation}
The coefficient of \eqref{eq:20.129} differs from the corresponding
coefficient in \eqref{eq:19.79} because we transform only the right-handed
chiral component of the electron field.  If we were to transform another
fermion field, of hypercharge $Y$, we would find a similar shift, with the
coefficient proportional to $Y^2$.  If this new field coupled to the $SU(2)$
or $SU(3)$ gauge fields, we would also find terms proportional to those field
strengths.  Thus, we can eliminate the term in \eqref{eq:20.127} involving
the $U(1)$ field strengths by making the change of variables \eqref{eq:20.128}
with $\alpha = -\theta_1$.  We can eliminate all three terms in
\eqref{eq:20.127} by making appropriate chiral rotations on three fermion
multiplets, say, $e_R$, $E_L$, and $Q_L$.  The change of variables
\eqref{eq:20.128}, which rotates the right-handed electron field, is not
symmetric under parity and, in fact, changes the definition of the parity
operation.  By making this change of variables, we are choosing new
coordinates in which the $P$ and $T$ transformation properties of the whole
theory are as simple as possible.

Let us now investigate the properties of the Lagrangian \eqref{eq:20.126}
under $P$, $C$, and $T$.  The couplings of the QCD gauge bosons are invariant
to each of these symmetries separately.  However, the couplings of the $SU(2)$
gauge bosons violate $P$ and $C$ as much as possible.  Recall from
Section~3.6 that $P$ converts a left-handed electron to a right-handed
electron, and that $C$ converts a left-handed electron to a left-handed
positron.  Each of these operations converts a particle that couples to
$SU(2)$ gauge bosons to one that does not.  However, the combination of these
two operations interchanges left-handed particles with right-handed
antiparticles.  Thus the combined operation $CP$ is a symmetry of
\eqref{eq:20.126}.  This Lagrangian is also invariant under time reversal.
Thus, we see that the discrete symmetries of $C$ and $P$, on the one hand, and
$CP$ and $T$, on the other, stand on a very different footing in gauge field
theories.  Any chiral gauge theory will naturally violate $C$ and $P$.  At
this level in our analysis, it is a mystery why $C$ and $P$ should be observed
to be approximate symmetries of Nature.  On the other hand, every theory of
gauge bosons and massless fermions respects $CP$ and $T$.  It is known
experimentally that Nature contains some interaction that violates $CP$, since
the $CP$ selection rules are weakly violated in the decays of the $K^0$ meson.
But to find a source for this violation, we must add terms to our basic gauge
theory \eqref{eq:20.126}.

We must, first of all, add dynamics to \eqref{eq:20.126} that will cause the
spontaneous breaking of $SU(2)\times U(1)$.  We will begin by working with the
simplest model with one Higgs scalar field $\phi$.  The most general
renormalizable Lagrangian for $\phi$ is
\begin{equation}
  \mathcal{L}_{\phi} = |D_\mu\phi|^2 + \mu^2\phi^{\dagger}\phi
    - \lambda(\phi^{\dagger}\phi)^2 .
  \label{eq:20.130}
\end{equation}
The Hermiticity of $\mathcal{L}_{\phi}$ implies that the parameters $\mu^2$
and $\lambda$ are real.  Thus this Lagrangian respects $P$, $C$, and $T$.  As
discussed at the end of the previous section, this Lagrangian also
automatically has the custodial $SU(2)$ symmetry required to produce the mass
relation \eqref{eq:20.117}.

Finally, we must add the terms that couple the Higgs field to the quarks and
leptons.  Here, renormalizability and gauge invariance provide the weakest
constraints, and there are many allowed interactions.  We will first analyze
the coupling of $\phi$ to the quark fields, and then generalize this
discussion to the leptons.

In writing the Higgs field couplings to the quarks, we should recall that
there are known to be three generations of quarks and leptons.  Thus there
are three doublets of left-handed quarks:
\begin{equation}
  Q^i_L = \begin{pmatrix}u_L\\ d_L\end{pmatrix},\qquad
  \begin{pmatrix}c_L\\ s_L\end{pmatrix},\qquad
  \begin{pmatrix}t_L\\ b_L\end{pmatrix} .
  \label{eq:20.131}
\end{equation}
There are six right-handed quarks, three with $Y = \frac{2}{3}$ and three
with $Y = -\frac{1}{3}$:
\begin{equation}
  u^i_R = (u_R, c_R, t_R),\qquad
  d^i_R = (d_R, s_R, b_R) .
  \label{eq:20.132}
\end{equation}
When we couple gauge fields to these quarks, we replace the ordinary
derivatives with covariant derivatives.  This automatically gives all of the
quarks the same coupling to QCD and all quarks of the same type the same
coupling to the weak and electromagnetic interactions.

The most general renormalizable Yukawa couplings of the Higgs field to the
quarks are
\begin{equation}
  \Delta\mathcal{L} = -\lambda^d_{ij}\ \bar{Q}^i_L\cdot\phi\,d^j_R
    - \lambda^u_{ij}\ \epsilon_{ab}\ \bar{Q}^{ia}_L\phi^{b*}\,u^j_R
    + \mathrm{h.c.} ,
  \label{eq:20.133}
\end{equation}
where $\lambda^d_{ij}$ and $\lambda^u_{ij}$ are arbitrary complex matrices in
generation space.  After spontaneous symmetry breaking, these terms give rise
to the quark mass matrices
\begin{equation}
  \bigl( M_d \bigr)_{ij} = \frac{v}{\sqrt{2}}\lambda^d_{ij},\qquad
  \bigl( M_u \bigr)_{ij} = \frac{v}{\sqrt{2}}\lambda^u_{ij} .
\end{equation}
These matrices are not in general diagonal.  We can diagonalize them by
making unitary transformations on the left- and right-handed quark fields:
\begin{equation}
  u_L = U^u_L u'_L,\qquad u_R = U^u_R u'_R,\qquad
  d_L = U^d_L d'_L,\qquad d_R = U^d_R d'_R .
  \label{eq:20.134}
\end{equation}
After these transformations, the quark mass matrices are diagonal, with
positive eigenvalues.  However, the charged-current weak interactions, which
couple $u_L$ to $d_L$ through the $W$ boson, are no longer diagonal in the
mass basis.  The mixing is described by the Cabibbo-Kobayashi-Maskawa (CKM)
matrix,
\begin{equation}
  V = U^{u\dagger}_L U^d_L .
  \label{eq:20.135}
\end{equation}
The CKM matrix is a unitary $3\times3$ matrix, which can be parametrized by
three angles and one $CP$-violating phase.  The presence of the phase means
that the CKM matrix provides a source of $CP$ violation in the weak
interactions, if the matrix is not real.  The existence of this phase is one
of the key predictions of the standard model, and it is the explanation of
the observed $CP$ violation in the $K^0$ system.

The same analysis applies to the leptons, if there is a right-handed
neutrino.  If the neutrinos are massless, there is no lepton mixing, and the
lepton analog of the CKM matrix is trivial.  If the neutrinos acquire
masses, through either a Dirac or a Majorana mass term, there is a
nontrivial lepton mixing matrix, the Pontecorvo-Maki-Nakagawa-Sakata (PMNS)
matrix.

\problems

\textbf{20.1} \emph{$SU(3)$ breaking by an adjoint scalar.}  Consider an
$SU(3)$ gauge theory with a scalar field in the adjoint representation, and
analyze the symmetry-breaking patterns for the vacuum expectation value
\eqref{eq:20.41}.  Compute the gauge boson masses explicitly, and verify the
results \eqref{eq:20.40} and \eqref{eq:20.42}.

\textbf{20.2} \emph{The Higgs mechanism in the Abelian case.}  Consider the
Abelian Higgs model of Section~20.1.  Work out the Feynman rules for the
theory in the unitarity gauge, and verify explicitly that the Goldstone boson
has been removed from the spectrum while the gauge boson has acquired a mass
and a longitudinal polarization.

\textbf{20.3} \emph{The $e^{+}e^{-}$ forward-backward asymmetry.}  Compute the
differential cross section for $e^{+}e^{-}\to f\bar{f}$ in the GWS theory,
including the effects of the $Z^0$ exchange diagram.  Show that, on the $Z^0$
resonance, the cross section has a forward-backward asymmetry whose
magnitude is determined by the couplings in \eqref{eq:20.80}.  Express the
asymmetry in terms of the parameter $\sin^2\theta_W$.

\textbf{20.4} \emph{Neutral-current neutrino scattering.}  Compute the
neutral-current contribution to deep inelastic neutrino scattering, using the
$Z^0$ couplings in \eqref{eq:20.80}.  Show that the ratio of the
neutral-current to charged-current cross sections is a measure of
$\sin^2\theta_W$, and compare to the experimental determination.

\textbf{20.5} \emph{The Goldstone boson equivalence theorem.}  In this
problem we derive the Goldstone boson equivalence theorem, which states that
the amplitude for emission of a longitudinal $W$ or $Z$ boson at high energy
is equal to the amplitude for emission of the corresponding Goldstone boson.
Show that, at energies much greater than $m_W$, the longitudinal
polarization vector of a massive gauge boson is approximately
\begin{equation}
  \epsilon^\mu_L(k) \approx \frac{k^\mu}{m_W} ,
\end{equation}
and that the coupling of this longitudinal mode is equivalent to the coupling
of the corresponding Goldstone boson.
